\chapter{Recipes}
\label{ch:recipes}

Checklists and code templates for every kind of change the tasks need. Each recipe names
the example that shows it working.

\section{Add a syscall (ex01)}
\label{sec:syscall}

\begin{center}
\footnotesize
\begin{tabular}{cL{6.6cm}L{8.2cm}}
\toprule
& \textbf{File} & \textbf{Add} \\
\midrule
1 & \file{common/include/kernel/syscall-definitions.h} & \code{\#define sc\_foo 1500} \\
2 & \file{common/include/kernel/Syscall.h} & \code{static size\_t foo(size\_t a, size\_t b);} \\
3 & \file{common/source/kernel/Syscall.cpp} & \code{case sc\_foo: return\_value = foo(arg1, arg2); break;} \\
4 & \file{common/source/kernel/Syscall.cpp} & the implementation \\
5 & \file{userspace/libc/include/nonstd.h} (or \code{pthread.h}, \ldots) & \code{extern int foo(int a, void* b);} \\
6 & \file{userspace/libc/src/nonstd.c} & \code{return \_\_syscall(sc\_foo, a, (size\_t) b, 0x00, 0x00, 0x00);} \\
\bottomrule
\end{tabular}
\end{center}

\begin{lstlisting}[style=cpp]
size_t Syscall::foo(size_t a, size_t user_ptr)
{
  if (user_ptr == 0 || user_ptr >= USER_BREAK - sizeof(Thing))  // check first
    return (size_t) -1;
  Thing copy = ...;                                            // work on kernel data
  *(Thing*) user_ptr = copy;                                    // one copy to user space at the end
  return 0;
}
\end{lstlisting}

\begin{trapbox}
\begin{itemize}
\item \code{return -1U;} is \code{0xFFFFFFFF} --- user space sees 4294967295. Use \code{(size\_t) -1}.
\item Forgotten \code{break;} $\to$ falls through into the next case.
\item \code{-Werror}: an unused parameter, a missing include, initializers out of
  declaration order (\code{-Werror=reorder}) all stop the build.
\item \file{nonstd.h} does not include \file{types.h} $\to$ \code{unknown type name 'size\_t'}.
\item At most 5 arguments. More: pass a pointer to a struct or array (A1 ``15-argument syscall'').
\item A new test \code{.c} needs cmake to run again (\code{run\_test.sh} does it) --- and a
  second \code{make}: a parallel build may copy the programs to the disk image before they
  are built (\code{loading /usr/shell.sweb failed}).
\end{itemize}
\end{trapbox}

\section{User pointers}
\label{sec:userptr}

\begin{enumerate}
\item \textbf{Check the whole range} below \code{USER\_BREAK}, overflow-safe:
  \code{p == 0 || p >= USER\_BREAK - size} $\to$ reject. (\code{p + size > USER\_BREAK} can wrap around.)
\item \textbf{Copy in, then check, then use.} Another thread of the process can change user
  memory between your check and your use (TOCTOU). Copy into a kernel variable first,
  validate the copy (A0-09: function pointers).
\item \textbf{Write results back at the end}, after all locks are released: a user page that
  is not mapped yet page-faults --- in kernel mode, on a user address. That is legal (it is
  loaded like a normal fault), but it must not happen while you hold a lock, and if no
  segment covers the address the process is killed right there (order your side effects:
  A0-08 ``touch the id pointer first'').
\item Never hand a user pointer to a driver or keep it for later: copy into kernel memory.
\item \textbf{Strings have no known length} --- there is no range to check up front. Copy byte by
  byte into a bounded kernel buffer, check \emph{each} address against \code{USER\_BREAK} before
  reading it, stop at the NUL or the limit, write the NUL yourself (\code{strncpy} does not when
  the source is long). Read user memory directly, not via \code{checkAddressValid}: that one
  calls a page ``unmapped'' that is merely not loaded yet (ex11, A0-11).
\end{enumerate}

\section{Physical pages and mappings (ex02, ex04)}
\label{sec:mapping}

\begin{center}
\footnotesize
\begin{tabular}{L{6.6cm}L{8.8cm}}
\toprule
\textbf{Call} & \textbf{Does} \\
\midrule
\code{PageManager::instance()->allocPPN()} & a free physical page, \textbf{zeroed}; asserts when memory is full \\
\code{PageManager::instance()->freePPN(ppn)} & gives it back (fills it with \code{0xFF}: use-after-free is detected) \\
\code{ArchMemory::getIdentAddressOfPPN(ppn)} & kernel address of that page --- valid in every address space \\
\code{arch\_memory\_.mapPage(vpn, ppn, user)} & maps \textbf{writable}; \code{user = 1}: user may access; \code{false} if vpn already mapped (caller keeps \code{ppn}); allocates page tables on the way \\
\code{arch\_memory\_.unmapPage(vpn)} & clears the entry, \textbf{frees the physical page}, frees empty page tables; asserts if not mapped; \textbf{no TLB flush} \\
\code{arch\_memory\_.resolveMapping(vpn)} & \code{ArchMemoryMapping m}: \code{m.pt[m.pti]} is the PTE, \code{m.page\_ppn}, \code{m.page\_size} (\code{PAGE\_SIZE} iff a 4 KiB page is mapped; \code{m.pt} is \code{nullptr} if there is no page table) \\
\code{arch\_memory\_.checkAddressValid(addr)} & ident address of the byte \code{addr} if mapped, else 0 --- write user memory without risking a fault \\
\code{ArchMemory::flushTlb()} & reload \code{cr3}: forget cached translations \\
\bottomrule
\end{tabular}
\end{center}

\code{vpn} and \code{ppn} are \textbf{page numbers} (\code{address / PAGE\_SIZE}), not addresses.
\textbf{One frame, one mapping:} \code{\textasciitilde ArchMemory} frees every present user page
at exit. A frame you map a second time (an alias, a shared page) must be unmapped
\emph{without} freeing first (\code{unmapPage(vpn, false)}), or it is freed twice: kernel panic
\code{"Double free PPN"} (ex12, A0-12, A1-10).

\begin{lstlisting}[style=cpp]
// map a fresh page (outside the lock: allocate; inside: map)
size_t ppn = PageManager::instance()->allocPPN();
loader->arch_memory_lock_.acquire();
bool mapped = loader->arch_memory_.mapPage(vpn, ppn, 1);
if (mapped && read_only)
{
  ArchMemoryMapping m = loader->arch_memory_.resolveMapping(vpn);
  m.pt[m.pti].writeable = 0;            // same critical section as mapPage
}
loader->arch_memory_lock_.release();
if (!mapped)
  PageManager::instance()->freePPN(ppn); // someone else mapped it first

// unmap
loader->arch_memory_lock_.acquire();
if (loader->arch_memory_.resolveMapping(vpn).page_size)
  loader->arch_memory_.unmapPage(vpn);   // frees the ppn - don't free it again
ArchMemory::flushTlb();
loader->arch_memory_lock_.release();
\end{lstlisting}

\textbf{TLB rule:} after \emph{removing} rights (unmap, read-only, NX, swap out) flush the TLB;
after \emph{adding} rights (map a non-present page, make writable) you need not --- a stale,
weaker entry just causes one more page fault that finds everything in order.

\textbf{The page-table entry} (\code{PageTableEntry}, \file{paging-definitions.h}):
\code{present, writeable, user\_access, write\_through, cache\_disabled, accessed, dirty,
size, global, ignored\_2 (3 bits), page\_ppn (28 bits), reserved\_1, ignored\_1 (11 bits),
execution\_disabled}. The CPU sets \code{accessed} on every access and \code{dirty} on every
write. When \code{present == 0}, the CPU ignores \emph{all} other bits --- you may store your
own information there (A2: the swap slot). \code{execution\_disabled} (NX) works: SWEB enables
\code{EFER.NXE} at boot.

\textbf{Walking all page tables} (A0-01, A2 ``all pages of a process''): copy the loop of
\code{ArchMemory::\textasciitilde ArchMemory()} --- four nested loops over the present
entries, PML4 index \code{0 .. 255} only (user half).

\section{A new page fault case (ex03)}
\label{sec:pagefault}

\begin{itemize}
\item \textbf{Where:} in \code{PageFaultHandler::handlePageFault}, inside
  \code{if (checkPageFaultIsValid(...))}, \emph{before} \code{loader\_->loadPage}. Only there
  you know the fault is a non-present user address.
\item \textbf{Flags} available: \code{address}, \code{user} (fault in user mode),
  \code{present} (page was present: protection violation), \code{writing},
  \code{fetch} (instruction fetch), \code{switch\_to\_us}.
\item Faults with \code{present == 1} never reach the valid branch: \code{checkPageFaultIsValid}
  rejects them. A protection case (copy-on-write, A2) needs a branch \emph{before} that check.
\item The ``kill'' exits: 9999 invalid fault (page fault handler), 666 no segment
  (\code{Loader::loadPage}), 999 file read failed, 888 CPU exception.
\item Returning without fixing the cause $\to$ the same fault again, forever.
\end{itemize}

\section{Changing the user registers (ex05)}
\label{sec:registers}

Allowed while handling a syscall or a page fault of that thread (its
\code{user\_registers\_} are saved and will be loaded on return), or for a thread that is
not scheduled (not yet added, or sleeping).

\begin{lstlisting}[style=cpp]
ArchThreadRegisters* regs = currentThread->user_registers_;
size_t rsp = ((regs->rsp - 128) & ~0xFULL) - sizeof(size_t); // red zone, align, ret slot
pointer slot = loader->arch_memory_.checkAddressValid(rsp);  // under the page-table lock
if (!slot) { /* no usable stack */ }
*(size_t*) slot = return_address;   // e.g. regs->rip (continue after the syscall) or 0
regs->rsp = rsp;
regs->rip = function;
regs->rdi = first_argument;         // rsi, rdx, rcx, r8, r9 for more
\end{lstlisting}

The syscall's own return value still goes into \code{rax} afterwards. Why the red zone matters:
\code{objdump -d prog.sweb} $\to$ \code{<\_\_syscall>} keeps its locals at \code{-0x10(\%rbp)
... -0x38(\%rbp)} while \code{rsp = rbp - 8} --- \emph{below} \code{rsp}.

\section{Threads (ex06, A0-08)}
\label{sec:threads}

\textbf{A user thread} = same \code{loader\_} (address space) + own user stack + own
registers + an entry in the scheduler:

\begin{lstlisting}[style=cpp]
class UserThread : public Thread { ... };          // Thread(..., Thread::USER_THREAD)
loader_ = process->loader_;                         // shared address space
/* map stack pages under the page-table lock */
ArchThreads::createUserRegisters(user_registers_, (void*) entry,
                                 (void*) (stack_top - sizeof(pointer)), getKernelStackStartPointer());
user_registers_->rdi = ...; user_registers_->rsi = ...;
ArchThreads::setAddressSpace(this, loader_->arch_memory_);
// later, when everything is set up:
Scheduler::instance()->addNewThread(thread);
\end{lstlisting}

\textbf{A kernel thread}: subclass \code{Thread} with \code{KERNEL\_THREAD}, implement
\code{Run()}, \code{new} it, \code{addNewThread}. When \code{Run()} returns, it is killed
and the CleanupThread deletes it --- never delete a thread yourself, never put one on
the stack.

\textbf{Ending a thread}: \code{currentThread->kill()} --- state \code{ToBeDestroyed}, never
returns. The object lives on until the CleanupThread deletes it: \textbf{destructors run in
the CleanupThread}, when the thread can never run again --- the right place to free its
stack. Another thread cannot be killed directly ``from outside'' safely: set a flag that it
checks (A1 \code{pthread\_cancel}).

\section{Waiting for something}
\label{sec:waiting}

\begin{lstlisting}[style=cpp]
// waiter                                   // the one who makes it true
lock.acquire();                             lock.acquire();
while (!condition_is_true)                  condition_is_true = true;
  condition_variable.wait();                condition_variable.signal();   // or broadcast()
/* use the state, still locked */           lock.release();
lock.release();
\end{lstlisting}

\begin{itemize}
\item \code{Condition(\&mutex, "name")}; \code{wait()} releases the mutex while sleeping and
  re-acquires it; SWEB \textbf{asserts that you hold no other lock} while waiting.
\item \code{signal()} wakes one waiter, \code{broadcast()} all (barriers, ``all threads done'').
\item \code{Scheduler::sleep()} / \code{wake(thread)} are the low-level primitives the locks are
  built on --- using them directly is easy to get wrong (lost wake-ups); prefer a Condition.
\item A \code{while (!x) yield();} loop is \emph{busy waiting}: acceptable for a quick test,
  not for a solution the tutor should accept (``join must sleep, not spin'').
\end{itemize}

\section{From another context into a process (ex07)}

Key presses are handled in the Console kernel thread, the timer in an interrupt --- neither
runs in the target process's context. Set a flag there (\code{ArchThreads::testSetLock(flag, 1)}
is an atomic exchange; \code{atomic\_add}, \code{atomic\_set} exist too) and let the target
do the work at a safe point in its own context (top of \code{syscallException}, page fault
handler, or a kernel thread you wake).

\section{Disk (ex08)}

\code{BDManager::getInstance()->getDeviceByName("idea2")} $\to$
\code{readData(offset, size, kernel\_buffer)} / \code{writeData(...)}: offset and size in
bytes, multiples of \code{getBlockSize()} (512); returns size or -1; \textbf{blocks} (yields)
until done --- interrupts on, no SpinLock held. \code{run\_test.sh} boots with a throw-away
copy of the disk (\code{-snapshot}).

\section{Kernel-wide objects}

\textbf{SWEB never runs constructors of global or static objects.} A global
\code{Mutex lock("x");} is zeroed memory and asserts on first use. Create such objects with
\code{new} in \code{startup()} (\file{main.cpp}) before the threads that use them exist, or
make them members of a singleton that is created at boot.

\section{Testing traps}
\label{sec:testing}

\begin{itemize}
\item \textbf{The timer runs at 18.2\,Hz} (\code{IRQ0\_TIMER\_FREQUENCY} is 0, the PIT keeps its
  default divisor 65536): one tick is 55\,ms. A round of \code{sched\_yield} passes the IdleThread,
  which halts until the next tick --- \textbf{300 yields take about 16 seconds}. Loops that yield
  look like hangs. Compute instead, or sleep on a Condition.
\item QEMU is fast: a busy loop of millions of iterations finishes within one 55\,ms time slice.
  To get a deterministic order between threads, let them yield a different number of times.
\item Lazy loading: a page of \code{.bss} or \code{.data} that was never touched has no frame yet
  (\code{vtop} $\to$ -1, \code{mprotect} fails, statistics are smaller than expected).
\item The disk image keeps every program ever built into it; after many tasks in one build folder
  \code{exe2minixfs} asserts (\code{name != ""}). Delete \file{SWEB-flat.vmdk} and \file{SWEB.qcow2}
  in the build folder (\code{run\_test.sh} and \code{task.sh} do it automatically).
\item Tests that must be \emph{killed} can only show it by the absence of a \code{[FAIL]} line and the
  kernel's message --- put them last in a test program.
\item Runner extras: \code{-n 5} (repeat: races), a final \code{exit} (SWEB's leak check at shutdown),
  \code{--keep-disk} (several boots on one disk), \code{'\&prog.sweb' '!sleep 1' '!key f3' help}
  (press a key while a program runs).
\item User-space libc is small: no \code{strcpy}, no \code{sprintf}, \code{malloc} is a stub;
  \code{exit} is not \code{noreturn} (GCC may warn about ``infinite recursion'').
\end{itemize}

\section{Debugging}

\begin{center}
\footnotesize
\begin{tabular}{L{7.4cm}L{8.0cm}}
\toprule
\textbf{Message} & \textbf{Meaning} \\
\midrule
\code{Maybe you are dereferencing a null-pointer} & fault below 4096 \\
\code{pagefault even though the address is mapped} & write to read-only / user access to a kernel-only page \\
\code{No section refers to the given address} & \code{loadPage}: address not in any ELF segment (exit 666) \\
\code{You are accessing a kernel address in user-mode} & user pointer $\geq$ \code{USER\_BREAK} \\
\code{Double free PPN} / \code{Page modified after free} & \code{freePPN} twice / write after free \\
\code{thread is holding unrelated locks while waiting on a condition variable} & \code{wait()} with another lock held \\
\code{name\_ is null, maybe lock was not initialised?} & a lock whose constructor never ran (global object!) \\
\code{don't use strcpy} & use \code{strncpy}/\code{memcpy} in the kernel \\
\code{You might be leaking physical memory pages} & at shutdown (\code{exit} in the shell): pages not freed \\
\code{Deadlock: Lock ... held already by currentThread} & you acquire a lock you already hold \\
\code{CIRCULAR DEADLOCK when waiting for ...} & lock-order cycle --- SWEB prints who holds what and waits for what \\
\code{Lock ... with IF=0 ... Now we're dead} & lock (or \code{new}) with interrupts disabled: interrupt handler, \code{schedule()} \\
\bottomrule
\end{tabular}
\end{center}

\begin{itemize}
\item \code{debug(FLAG, "fmt", ...)}: categories in \file{debug.h}; add your own
  (\code{const size\_t MYFLAG = Ansi\_Cyan | OUTPUT\_ENABLED;}). \code{LOADER},
  \code{PAGEFAULT}, \code{SYSCALL} are on by default.
\item In SWEB: \textbf{F9} page bitmap + kernel memory, \textbf{F10} locks, \textbf{F11} stack
  traces of all threads, \textbf{F12} thread list.
\item \code{run\_test.sh --raw} shows the whole kernel log; \code{-n 10} repeats a run.
\item gdb: \code{make debug \&\& make -j \&\& make qemugdb} in the build folder, then
  \code{gdb} in a second terminal (see \file{common\_commands.md}).
\end{itemize}
